\def\level{BTS}  % Classe à droite
\def\course{CIEL 2}
\def\chaptername{Lois continues} % Titre au centre
\def\docpart{Loi normale~--~Exercices} % Partie du cours
\def\docname{C2252}        % Identifiant à gauche

\pagestyle{cours_FP}
\makeheaderBTS{} % Affiche le bandeau

\begin{multicols}{2}
\begin{exercice_cours}[\quad]{}


\includegraphics[width=\linewidth]{01.png}

Associer chaque courbe à la densité de probabilité qu'elle représente.


\includegraphics[width=\linewidth]{03.png}


\end{exercice_cours}
\begin{exercice_cours}[\quad]{}


\includegraphics[width=\linewidth]{02.png}

Associer chaque courbe à la densité de probabilité qu'elle représente.

\includegraphics[width=\linewidth]{04.png}


\end{exercice_cours}
\begin{exercice_cours}[\quad]{}


Soit une variable aléatoire $X$ suivant la loi normale centrée réduite $\mathscr{N}(0;1)$

\begin{itemize}	\item $P(1<X<2)=$\dotfill
	\item $P(-0,5<X<1,3)=$\dotfill
	\item $P(X>-0,02)=$\dotfill
	\item $P(X<1,28)=$\dotfill
\end{itemize}
\end{exercice_cours}
\begin{exercice_cours}[\quad]{}


Soit une variable aléatoire $X$ suivant la loi normale centrée réduite $\mathscr{N}(15;2)$

\begin{itemize}	
	\item $P(10<X<20)=$\dotfill
	\item $P(X>16)=$\dotfill
	\item $P(X<18)=$\dotfill
\end{itemize}

\end{exercice_cours}
\begin{exercice_cours}[\quad]{}


Le délai de livraison $X$ d’une pièce, en jours, suit la loi normale d’espérance $30$ et d’écart-type $5$.
Quelle est la probabilité pour que le délai de livraison soit :

\begin{enumerate}
	\item Compris entre 22 et 38 jours ? \dotfill
	\item Inférieur à 27 jours ? \dotfill
	\item Supérieur à 35 jours ? \dotfill
\end{enumerate}




\end{exercice_cours}
\begin{exercice_cours}[\quad]{}


Une entreprise fabrique entre autres des câbles RJ45 croisés.

On désigne par $X$ la variable aléatoire qui, à tout câble tiré au hasard de la production d'une journée associe sa longueur, exprimée en centimètres.

On admet que $X$ suit la loi normale de moyenne $500$ et d'écart type $2$.

On prélève au hasard un câble RJ45 croisé dans la production d'une journée.

\begin{enumerate}
	\item Calculer la probabilité de l'événement A :
	
	"La longueur du câble RJ45 croisé est inférieure à 501,9 centimètres". \dotfill
	\item Calculer la probabilité de chacun des événements suivants :
	\begin{enumerate}
		\item La longueur du câble RJ45 croisé est supérieure à 501,6 cm. \dotfill
		\item La longueur du câble est comprise entre 501,6 et 501,9 cm. \dotfill
	\end{enumerate}
\end{enumerate}

\end{exercice_cours}
\begin{exercice_cours}[\quad]{}


Une machine fabrique des gaines d'isolation utilisées dans les réseaux informatiques. Le diamètre attendu pour ces gaines est de $12$ mm.

Cependant, le diamètre réel $X$ (en mm) fluctue légèrement de manière aléatoire autour de cette valeur : il suit la loi normale d'espérance $12$ et d'écart-type $0,08$.

Les gaines de diamètre inférieur à $11,9$ ou supérieur à $12,2$ sont rejetées.

Calculer la probabilité qu'une gaine prise au hasard soit rejetée lors du test de mesure du diamètre ?


\end{exercice_cours}
\begin{exercice_cours}[\quad]{}


Soit $X$ une v.a. qui suit une loi normale d'espérance $6$.

On sait que $P(5,5\geq X\geq6,5)= 0,95$.

Déterminer $\sigma$ à $10^{-3}$ près : \dotfill


\end{exercice_cours}


\end{multicols}	